\section{从分子到细胞：跨尺度的 AI 生物学研究}

\subsection{超越蛋白质设计的更广阔图景}

Amini 强调，蛋白质设计只是生物学 AI 这个更大图景中的一个层次。生物学是一个\textbf{深度层级化的系统}：

\begin{table}[H]
\centering
\caption{生物学的层级结构与对应的 AI 挑战}
\begin{tabular}{@{}llll@{}}
\toprule
\textbf{层级} & \textbf{对象} & \textbf{数据类型} & \textbf{AI 任务} \\
\midrule
分子层 & 蛋白质、DNA & 序列、结构 & 分子设计、功能预测 \\
细胞层 & 单细胞状态 & 转录组、蛋白组 & 细胞状态表征、扰动预测 \\
组织层 & 组织切片 & 病理图像 & 疾病检测、分类、预后 \\
个体层 & 患者数据 & 多组学、临床记录 & 精准医疗、疗效预测 \\
\bottomrule
\end{tabular}
\end{table}

MSR BioML 团队的研究跨越了所有这些层级，致力于开发能在不同尺度上推理的 AI 方法，最终将它们整合起来驱动对临床有影响力的决策。

\subsection{序列与结构的融合方向}

在蛋白质设计领域，团队正在积极探索序列与结构信息的\textbf{联合建模}。Amini 指出，序列和结构是互补的：

\begin{itemize}
    \item \textbf{结构}提供高分辨率的精确几何控制
    \item \textbf{序列}提供更广泛的生物多样性覆盖
\end{itemize}

将两者结合的途径包括：将序列和结构映射到共同的表征空间、联合训练序列-结构生成模型等。

\subsection{文本-蛋白质多模态}

另一个前沿方向是蛋白质与自然语言的\textbf{多模态对齐}：

\begin{itemize}
    \item 给定蛋白质序列/结构 $\rightarrow$ 生成功能描述文本
    \item 给定功能描述文本 $\rightarrow$ 设计满足描述的蛋白质
\end{itemize}

这朝着 Amini 在讲座开头描绘的愿景迈进：用自然语言提示（如"设计一个能靶向并结合乳腺癌细胞的蛋白质"）来驱动蛋白质设计。

\subsection{细胞尺度的 AI：理解疾病状态}

在细胞层面，团队正在开发能够定义和表征\textbf{细胞状态} (cellular states) 的 AI 方法。目标是理解疾病（特别是癌症）如何改变细胞状态，并利用这些信息来指导治疗策略的选择。

\subsection{数字病理学：AI 驱动的组织图像分析}

在组织层面，MSR 团队在\textbf{数字病理学} (digital pathology) 领域取得了突破性进展——将计算机视觉方法应用于组织病理图像的分析，实现癌症的自动检测、分类和预后判断。

\begin{knowledgebox}{MSR 与 Broad Institute 的战略合作}
MSR BioML 团队与 MIT 和 Harvard 的 Broad Institute 建立了长期战略研究合作。Broad Institute 是全球精准医学和基因组学研究的领军机构。双方的合作目标是将 AI 能力与顶尖的实验和临床资源紧密结合，实现从患者分子数据到 AI 预测再到实验验证再到临床应用的完整闭环。这种合作模式体现了 AI for Biology 领域中计算与实验缺一不可的核心特征。
\end{knowledgebox}

\subsection{本章小结}

AI for Biology 的研究远不止蛋白质设计，而是跨越分子、细胞、组织、个体多个层级的系统性工程。MSR BioML 团队的工作覆盖从蛋白质序列/结构联合建模到数字病理学，最终目标是将多尺度的 AI 方法整合起来，驱动对临床实践有实质影响的精准医疗决策。


